\IfFileExists{../main.tex}{%
  \documentclass[../main.tex]{subfiles}%
}{%
  \documentclass[main.tex]{subfiles}%
}
\begin{document}
\setcounter{chapter}{5}
\makeatletter
\newcommand{\chapterFiveRef}[2]{%
  \@ifundefined{r@#1}{#2}{\ref{#1}}%
}
\makeatother

% ===========================================================================
\section{Tổng quan}

Chương này trình bày những giải pháp kỹ thuật nổi bật nhất được phát triển trong quá trình thực hiện đồ án. Đây phần lớn là các bài toán nảy sinh từ đặc thù nghiệp vụ chia sẻ thực phẩm. nhiều bên tham gia, thao tác đồng thời, quy trình kéo dài và dễ bị bỏ ngang, dữ liệu nhạy cảm về vị trí và danh tính. Mỗi mục dưới đây gồm ba phần: bài toán, giải pháp và kết quả; đồng thời chỉ ra khả năng tổng quát hoá sang lớp bài toán rộng hơn khi phù hợp. Các nội dung đã mô tả ở chương trước chỉ được nhắc lại vắn tắt kèm tham chiếu chéo.

% ===========================================================================
\section{Ghép cặp an toàn trong môi trường tương tranh}
\label{section:5.1}

\textbf{Bài toán.}

Mỗi đơn quyên góp là một tài nguyên độc quyền: tại một thời điểm chỉ được phép gắn với một người nhận hoặc một tình nguyện viên thực hiện giao nhận. Tuy nhiên, hệ thống vận hành theo nguyên tắc đến trước được phục vụ trước (đã trình bày tại mục~\chapterFiveRef{subsection:2.3.6}{2.3.6}), do đó nhiều người dùng có thể gửi yêu cầu ``Kết nối'' hoặc ``Nhận đơn'' gần như đồng thời.

Nếu áp dụng cách xử lý tuần tự theo mô hình ``kiểm tra rồi cập nhật'' (check-then-set), hai yêu cầu song song có thể cùng quan sát thấy đơn vẫn còn khả dụng và đều thực hiện cập nhật thành công. Hệ quả là xuất hiện hiện tượng double-booking, trong đó nhiều người dùng cùng cho rằng mình đã nhận được một đơn quyên góp, hoặc hệ thống tạo ra nhiều bản ghi giao nhận (Delivery) cho cùng một đơn.

\textbf{Giải pháp.}

Để bảo đảm tính nhất quán dữ liệu trong môi trường có truy cập đồng thời, hệ thống sử dụng cơ chế cập nhật nguyên tử của cơ sở dữ liệu để hiện thực mẫu compare-and-set. Thay vì tách riêng bước kiểm tra và bước cập nhật, mỗi thao tác ``chốt'' đơn được thực hiện thông qua một câu lệnh cập nhật có điều kiện, chỉ thành công khi đơn vẫn đang ở trạng thái hợp lệ và chưa được người khác nhận.

Nhờ đặc tính nguyên tử của thao tác ghi, các yêu cầu đồng thời sẽ được cơ sở dữ liệu xử lý theo cơ chế tuần tự hóa. Kết quả là chỉ một yêu cầu duy nhất thỏa mãn điều kiện cập nhật và chiếm được tài nguyên, trong khi các yêu cầu còn lại thất bại do điều kiện không còn đúng. Cách tiếp cận này loại bỏ nguy cơ double-booking mà không cần sử dụng khóa tường minh hoặc các giao dịch nhiều bước có chi phí cao.

\begin{lstlisting}[morekeywords={const,await,async,throw,new,null},keywordstyle=\bfseries,frame=single,framesep=4pt,xleftmargin=10pt,escapeinside={(@}{@)}]
const claim = await FoodDonation.updateOne(
  { _id, status:'PENDING',selected_receiver_id: null },  
  (@\textit{// điều kiện}@)
  { $set: { selected_receiver_id, selected_at } }          
  (@\textit{// hành động}@)
);
if (claim.modifiedCount === 0) throw Error(409);           
(@\textit{// thua cuộc đua}@)
\end{lstlisting}

Tầng lưu trữ dữ liệu bảo đảm các thao tác cập nhật đồng thời được thực hiện theo cơ chế tuần tự hóa. Do đó, trong trường hợp nhiều yêu cầu cùng tác động lên một tài nguyên, chỉ một yêu cầu duy nhất thỏa mãn điều kiện cập nhật và nhận được kết quả modifiedCount = 1, trong khi các yêu cầu còn lại nhận modifiedCount = 0 và được hệ thống phản hồi bằng mã lỗi HTTP 409 (Conflict). Cơ chế này giúp ngăn ngừa các tình huống cập nhật cạnh tranh và bảo đảm tính nhất quán của dữ liệu.

Nguyên tắc tương tự được áp dụng cho nghiệp vụ tình nguyện viên nhận đơn quyên góp thông qua thao tác findOneAndUpdate với điều kiện bảo vệ (guard) status = PENDING, cũng như cho quá trình tạo bản ghi giao nhận (Delivery).

Đối với Delivery, do hệ thống cần tạo bản ghi mới trước khi liên kết với đơn quyên góp, luận văn áp dụng quy trình tạo -- chiếm quyền -- hoàn tác (create--claim--rollback). Cụ thể, hệ thống trước hết tạo một bản ghi Delivery, sau đó thực hiện thao tác chiếm quyền (claim) một cách nguyên tử trên điều kiện delivery\_id = null. Nếu thao tác chiếm quyền không thành công do đã có yêu cầu khác thực hiện trước, bản ghi Delivery vừa tạo sẽ được xóa bỏ nhằm tránh phát sinh các dữ liệu không còn được tham chiếu trong hệ thống.


\textbf{Tổng quát hoá.} 

Đây chính là lời giải kinh điển cho lớp bài toán cấp phát tài nguyên hữu hạn dưới tương tranh (đặt vé, giữ chỗ, rút tồn kho\ldots): thay vì khoá bi quan, ta đặt điều kiện bất biến vào ngay mệnh đề ghi và để tầng lưu trữ làm trọng tài.

\textbf{Kết quả.} 

Kịch bản hai người nhận bấm kết nối đồng thời luôn cho đúng một người thắng, người còn lại nhận thông báo ``Đơn vừa được người khác nhận'' (ca kiểm thử TC09, mục~\chapterFiveRef{section:4.4}{4.5}); không quan sát thấy đơn bị chốt đôi hay Delivery trùng trong suốt quá trình thử nghiệm.

% ===========================================================================
\section{Ước lượng khoảng cách lai: nhanh, tiết kiệm và bền bỉ}
\label{section:5.2}

\textbf{Bài toán.} 

Trải nghiệm cốt lõi của ứng dụng là ``đơn ở gần hiển thị trước''. Khoảng cách đường chim bay (Haversine, đã trình bày ở Chương 3, công thức~(\chapterFiveRef{eq:haversine}{3.1})) tính rất nhanh nhưng sai lệch so với quãng đường di chuyển thực tế. Ngược lại, khoảng cách đường bộ từ dịch vụ bản đồ thì chính xác nhưng (i) tốn hạn ngạch (quota) và tiền theo từng lượt gọi, (ii) có độ trễ mạng và (iii) có thể thất bại. Ngoài ra dữ liệu toạ độ thực tế không sạch: hồ sơ mới khởi tạo mặc định (0,0) --- rơi vào Vịnh Guinea sẽ cho khoảng cách sai hàng nghìn km.

\textbf{Giải pháp.} 

Tác giả thiết kế một bộ ước lượng khoảng cách hai tầng có suy giảm an toàn, gồm bốn thành phần. (i)~Tầng chính gọi Google Distance Matrix để lấy quãng đường đường bộ khi có khoá API. (ii)~Tầng dự phòng: nếu không có khoá, API lỗi hoặc mất mạng, hệ thống tự động quay về Haversine, nên người dùng vẫn luôn có thứ tự sắp xếp hợp lý thay vì màn hình trống. (iii)~Bộ nhớ đệm (cache) thông minh: kết quả đường bộ được lưu in-memory với thời gian sống 10 phút, khoá cache là cặp toạ độ đã làm tròn tới ~111\,m; nhờ làm tròn, nhiều người nhận khác nhau ở gần nhau cùng xem một danh sách sẽ dùng chung một mục cache, giảm mạnh số lượt gọi API, đồng thời khi người quyên góp đổi địa chỉ thì khoá cache đổi theo nên tự hết hiệu lực. (iv)~Vệ sinh dữ liệu: mọi toạ độ (0,0) hoặc ngoài miền hợp lệ bị loại trước khi tính.
Bán kính cũng được dùng để lọc đối tượng nhận thông báo (10\,km cho người quyên góp khi có yêu cầu mới, 20\,km cho tình nguyện viên khi điều phối) nhằm tránh làm phiền người ở quá xa.

\textbf{Tổng quát hoá.} 

Mẫu ``nguồn chính bên ngoài + dự phòng nội bộ + cache khử trùng lặp'' áp dụng được cho mọi tích hợp dịch vụ trả phí/không ổn định (dịch máy, OCR, gợi ý\ldots): ưu tiên chất lượng cao khi có thể, nhưng không bao giờ để toàn hệ thống phụ thuộc sống còn vào bên thứ ba.

\textbf{Kết quả.} 

Danh sách đơn luôn sắp xếp được theo khoảng cách kể cả khi không cấu hình khoá bản đồ; khi có khoá, cache giúp cắt giảm phần lớn lượt gọi lặp lại trong cùng khu vực; loại bỏ hoàn toàn lỗi ``khoảng cách hàng nghìn km'' do toạ độ rác.

% ===========================================================================
\section{Vòng đời đơn như một máy trạng thái tự phục hồi}
\label{section:5.3}

\textbf{Bài toán.} 

Một đơn đi qua nhiều bước và nhiều bên (người quyên góp, người nhận, tình nguyện viên). Trong thực tế, người dùng bỏ ngang là chuyện thường: người nhận chọn ``tự lấy'' rồi không bao giờ đến; tình nguyện viên báo ``đang giao'' rồi biến mất; người nhận nhận hàng nhưng quên bấm xác nhận. Nếu không xử lý, các đơn này kẹt trạng thái vĩnh viễn, làm sai lệch thống kê và giam giữ tài nguyên (đơn không bao giờ được mở lại cho người khác, điểm thưởng không bao giờ được trao).

\textbf{Giải pháp.} 

Tác giả mô hình hoá toàn bộ vòng đời đơn thành một máy trạng thái tường minh (Hình~\ref{fig:state_lifecycle}), rồi bổ sung ba tác vụ nền định kỳ đóng vai trò ``người dọn dẹp'' (janitor), tự đưa các trạng thái treo về trạng thái kết thúc. (i)~expireOverdueDonations (mỗi 15 phút) đưa đơn PENDING quá hạn sử dụng sang EXPIRED, đồng thời mở lại yêu cầu thực phẩm liên kết (nếu có) để người quyên góp khác còn cơ hội đáp ứng. (ii)~autoConfirmStaleDeliveries (mỗi 1 giờ) xử lý đơn ở AWAITING\_CONFIRMATION quá 24 giờ mà người nhận chưa bấm xác nhận: tự xác nhận hoàn tất và trao điểm. (iii)~cancelStaleOnTheWayDeliveries (mỗi 30 phút) đưa đơn ở ON\_THE\_WAY quá 6 giờ (tình nguyện viên ``no-show'') sang CANCELLED.
Điểm mấu chốt là mỗi lượt quét đều \textbf{bất biến với lần chạy lại} (idempotent): mọi chuyển trạng thái dùng cùng cơ chế cập nhật có điều kiện ở mục~\ref{section:5.1}, nên dù tác vụ chạy chồng lấn hay đơn vừa được người dùng thao tác xen vào, cũng không gây chuyển trạng thái sai.

\begin{figure}[H]
    \centering
    \chaptergraphic[width=0.95\textwidth]{Hinhve/Giaiphap/state_lifecycle.png}
    \caption{Máy trạng thái vòng đời đơn với các chuyển trạng thái tự phục hồi (cron)}
    \label{fig:state_lifecycle}
\end{figure}

\textbf{Tổng quát hoá.} 

Đây là mẫu đối soát theo thời hạn (timeout-driven reconciliation) cho mọi quy trình dài có khả năng bị bỏ ngang: thay vì tin rằng người dùng luôn hoàn thành quy trình, hệ thống định kỳ tự đưa trạng thái về nhất quán.

\textbf{Kết quả.} 

Không còn đơn ``kẹt'' vô thời hạn; quyền lợi được bảo đảm công bằng kể cả khi một bên quên thao tác (ví dụ người nhận quên xác nhận thì sau 24 giờ người quyên góp và tình nguyện viên vẫn được cộng điểm).

\section{Tổng kết chương}
\label{section:5.7}

Các giải pháp trên tuy thuộc những khía cạnh khác nhau nhưng cùng chia sẻ một triết lý thiết kế: \ưu tiên tính đúng đắn dưới điều kiện không lý tưởng, người dùng thao tác đồng thời, dịch vụ ngoài có thể lỗi, quy trình có thể bị bỏ ngang, và luôn tồn tại động cơ gian lận. Nhiều mẫu trong số này (compare-and-set, dự phòng + cache, đối soát theo thời hạn, bộ đếm phiên bản) có thể tái sử dụng cho nhiều hệ thống điều phối nguồn lực khác ngoài phạm vi chia sẻ thực phẩm.

\ifSubfilesClassLoaded{\printbibliography{}}{}
\end{document}
